% =======================================================================
\chapter*{Siglas y nomenclatura}
\addcontentsline{toc}{chapter}{Siglas y nomenclatura}
% =======================================================================

\section*{Siglas}

\begin{tabular}{lp{11cm}}
\toprule
\textbf{Sigla} & \textbf{Significado} \\
\midrule
ACO   & \emph{Ant Colony Optimization} — optimización por colonias de hormigas \\
AMPL  & \emph{A Mathematical Programming Language} — lenguaje de modelado para
        optimización matemática \\
BKS   & \emph{Best Known Solution} — mejor valor conocido para una instancia,
        reportado en la literatura \\
CLI   & \emph{Command Line Interface} — interfaz de línea de comandos \\
CSV   & \emph{Comma-Separated Values} — formato tabular de texto plano \\
CVRP  & \emph{Capacitated Vehicle Routing Problem} — ruteo de vehículos con
        capacidad \\
GA    & \emph{Genetic Algorithm} — algoritmo genético \\
GAP   & Distancia porcentual entre la solución obtenida y el BKS \\
JSON  & \emph{JavaScript Object Notation} — formato de intercambio de datos \\
JSSP  & \emph{Job Shop Scheduling Problem} — planificación en taller de trabajos \\
MILP  & \emph{Mixed Integer Linear Programming} — programación lineal entera mixta \\
NEH   & Heurística constructiva de Nawaz, Enscore y Ham para Flowshop \\
\NP   & Clase de complejidad computacional \emph{Nondeterministic Polynomial time} \\
PFSP  & \emph{Permutation Flowshop Scheduling Problem} — planificación de flujo
        en taller \\
POC   & Problema de Optimización Combinatoria \\
SA    & \emph{Simulated Annealing} — recocido simulado \\
TS    & \emph{Tabu Search} — Búsqueda Tabú \\
TSP   & \emph{Travelling Salesman Problem} — problema del viajante \\
TSPLIB& Biblioteca pública de instancias de referencia para TSP y variantes \\
\bottomrule
\end{tabular}

\vspace{2em}

\section*{Nomenclatura}

\begin{tabular}{lp{11cm}}
\toprule
\textbf{Símbolo} & \textbf{Significado} \\
\midrule
$f^*$            & Mejor valor objetivo encontrado por el algoritmo \\
$\mathcal{T}$    & Lista tabú \\
$|\mathcal{T}|$  & Tamaño de la lista tabú \\
$T_{\max}$       & Número máximo de iteraciones \\
$L$              & Umbral de estancamiento que dispara el reinicio \\
$\lambda$        & Coeficiente de penalización adaptativa por frecuencia \\
$\Phi$           & Memoria de frecuencia de movimientos \\
$C_{\max}$       & \emph{Makespan}: instante de término del último trabajo \\
$\sigma$         & Desviación estándar sobre las semillas \\
pp               & Puntos porcentuales \\
\bottomrule
\end{tabular}

\vspace{1.5em}

\noindent Los tres niveles de Búsqueda Tabú se denotan a lo largo del documento
como \textbf{TS-Básico}, \textbf{TS-Intermedio} y \textbf{TS-Avanzado}.
